% Separate post-CP24 candidate. All statements await full independent audits.
\subsection{Irreducible residual inertia at arbitrary prime-to-\(p\) ramification}

Throughout this subsection \(p\ge5\), \(q=p^2\), \(N=q-1\), and
\(K\) is complete discretely valued of mixed characteristic \((0,p)\)
with algebraically closed residue field \(k\), \(\mu_p\subset K\),
and absolute ramification \(e\) prime to \(p\). The two height-two
groups \(M_1,M_2\) acquire one-dimensional formal good reduction,
and their actual \(p^2\)-torsion is isomorphic. Their common residual
inertia representation is irreducible over \(\mathbf F_p\).
We prove equality of their finite special inertia characters.
No restriction \(e=p-1\), CM hypothesis, or change of Tate lattice
is imposed.

\begin{lemma}[Common fields, coordinates, and the new normalization]
\label{lem:B-RD67-normalization}
\status{Candidate; complete arbitrary-index derivation written.}
Take the tame extension \(L/K\) of degree \(N\), whether or not
it acts faithfully on either special fibre. Both actual groups
have good formal models over \(L\), and all their \(p\)-torsion
is \(L\)-rational. Their common second-level division field \(F/L\)
has elementary abelian group
\[
 W\subset M_2(\mathbf F_p),\qquad |W|=p^r,\quad r\in\{2,3,4\}.
 \tag{B-RD67.1}
\]
Here \(B\in W\) acts on the torsion module by \(1+pB\).
There are a tame complement \(\langle\sigma\rangle\) of order \(N\),
an exact eigen-uniformizer \(\varpi\) of \(F\), and integral
Honda-framed coordinates \(T_i\) such that
\[
 \begin{gathered}
 \sigma\varpi=\eta\varpi,\quad \eta^{p^r}=\zeta,\quad
 T_i(\sigma P)=\zeta^{c_i}\sigma(T_i(P)),\\
 A=v_F(T_i(P))=e p^{r-2}\quad(P\text{ primitive of order }p^2),\\
 v_F(T_i(Q))=qA\quad(0\ne Q\in M_i[p]),\qquad
 e_F=NqA .
 \end{gathered}
 \tag{B-RD67.2}
\]
The root \(\zeta\) is primitive of order \(N\), defined by the action
of \(\sigma\) on a tame uniformizer of \(L\). Put \(b_0=Ne\).
Primitive leading coefficients identify the two residual spaces
with \(d_i\mathbf F_q\subset k\). Their tame eigenvalue pairs are
\[
 \{\zeta^{c_i+e},\zeta^{p(c_i+e)}\}.
 \tag{B-RD67.3}
\]
If the special character pairs \(\{c_i,pc_i\}\) differ modulo \(N\),
the residual intertwiner in these coordinates is
\(z\mapsto a z^p\), with \(a\in\mathbf F_q^\times\).
\end{lemma}
\begin{proof}
Lemma~\ref{lem:B-RD49-inertia} gives actual good reduction over
a tame degree dividing \(N\). A normal \(p\)-subgroup of an
irreducible two-dimensional \(\mathbf F_p\)-representation acts
trivially: its nonzero fixed subspace is stable under the full
group. The remaining cyclic irreducible residual image lies in
\(\mathbf F_q^\times\), so \(L\) also contains the residual division
field. The second-level kernel is the additive matrix space.
After adjoining \(L\) there is no further prime-to-\(p\) quotient
in that kernel. Thus \(F/L\) is the indicated elementary abelian
extension.

The valuation filtration on \(M_i[p]\) is inertia stable.
Irreducibility makes all nonzero first-level coordinates have
one valuation. The product of the \(q-1\) nonzero torsion
coordinates, or equivalently the linear coefficient of the
Weierstrass polynomial of \([p]\), gives that valuation
\(e_L/N=e\). The one-segment Newton polygon has endpoints
\((1,eN),(q,0)\). Subtracting the coordinate of a nonzero
first-level point gives one inverse segment from \((0,e)\) to
\((q,0)\): all intermediate coefficients are strictly above it.
Every primitive second-level coordinate therefore has
\(v_L\)-valuation \(e/q\). Its integral valuation in \(F\) is
\(e p^{r-2}\); since \(p\nmid e\), integrality forces \(r\ge2\).
This proves all valuations in (B-RD67.2).

The extension of \(C_N\) by the additive group \(W\) splits.
Indeed, if a lift \(s\) of a generator has \(s^N=w\), then \(w\)
is fixed by \(s\); replacing \(s\) by \(s(-N^{-1}w)\) makes its
\(N\)-th power one. Averaging a uniformizer against its leading
character gives the exact eigen-uniformizer. Its eigenvalue
satisfies \(\eta^{p^r}=\zeta\), hence is primitive of order \(N\).

Tate full faithfulness extends the actual semilinear descent maps
and their inverses to the good models
\cite[Theorem 4 and Corollaries 1--2, p.~181]{Tate1967PDivisible}.
The finite special action can be put into a standard Honda
coordinate integrally. Explicitly, for a special unit \(v\) of
order \(m\mid N\) with residue equal to a standard root \(\lambda\),
\[
 x=\sum_{j=0}^{m-1}v^j\lambda^{-j}
 \quad\text{satisfies}\quad vx=x\lambda,\qquad
 x\equiv m\pmod\Pi .
\]
Thus \(x\) is a unit in the height-two maximal order, not merely
a rational conjugator. Averaging an integral lift of the resulting
Honda parameter gives the displayed exact semilinear coordinate;
its special \(p\)-series remains \(T^q\).

Write \(R_i=M_i[p^2](F)\), \(V_i=R_i/pR_i\). The leading
coefficient of \(T_i(P)/\varpi^A\) gives an injective
\(\mathbf F_p\)-linear map \(V_i\to k\): two primitive points
have the same leading coefficient exactly when their difference
is first-level. Multiplication by \(p\) raises this coefficient
to its \(q\)-th power, since the degree-\(q\) term of \([p]\)
is uniquely lowest at a primitive point. The first-level leading
space, acted on irreducibly by tame inertia, is a scalar multiple
of \(\mathbf F_q\), and so is the primitive leading space.
The uniformizer contribution is
\(\eta^A=\zeta^{p^{-r}e p^{r-2}}=\zeta^e\).
This proves (B-RD67.3). Its eigenvalues are not in \(\mathbf F_p\).
An intertwiner between the two multiplication actions is either
\(\mathbf F_q\)-linear or Frobenius-semilinear. The linear case
forces \(c_1=c_2\pmod N\), giving equal special characters.
Thus different special characters require the latter case.
Under multiplication by \(p\), normalized leading coordinates
are \(z\mapsto z^q=z\); the same intertwiner acts on the
first-level submodule. This fixes the orientation used below.
\end{proof}

\begin{lemma}[Exact coefficient calculus and the wild images of dimensions two and three]
\label{lem:B-RD67-small-wild}
\status{Candidate; all lower-break cases written explicitly.}
Under the preceding hypotheses, a different-special-character
comparison cannot have \(r=2\) or \(r=3\).
\end{lemma}
\begin{proof}
Choose a coefficient ring in \(\mathcal O_K\), and representatives
of every residue class in that ring, by
\cite[Theorem 10.160.8 and proof]{Stacks032A}.
They are fixed by \(W\). Since \(e_F=NqA>qA\), all calculations
modulo \(\varpi^{qA+1}\) with these representatives are
characteristic-\(p\) calculations. For \(g\ne1\) define its lower
break and leading coefficient by
\[
 v_F(g\varpi-\varpi)=b_g+1,\qquad
 d_g=\overline{(g\varpi-\varpi)/\varpi^{b_g+1}} .
 \tag{B-RD67.4}
\]
A monomial of degree \(j\), with \(p^s\Vert j\), first changes
in degree \(j+p^s b_g\), with coefficient
\((j/p^s)d_g^{p^s}\) modulo \(p\). Terms divisible by \(p\)
in the mixed-characteristic expansion occur beyond \(qA\).
On a subgroup having one common break, the leading map
\(g\mapsto d_g\) is additive and injective. Tame conjugation
acts on it by \(\eta^{b_g}\).

If \(r=2\), then \(A=e\) is prime to \(p\). Every primitive
coordinate changes first in degree \(A+b_g\), with nonzero
coefficient \(\bar e\,d_i z d_g\). Thus every \(g\ne1\)
displaces every primitive point, and \(A+b_g=qA\). The matrix
\(B_i(g)\) on the residual space consequently acts by
\[
 B_i(g)z=\kappa_i d_g z,\qquad \kappa_i\ne0.
 \tag{B-RD67.5}
\]
Its image as \(g\) varies is the entire Cartan field
\(\mathbf F_q\), since \(W\) has order \(q\).
The different-character Frobenius intertwiner instead requires
the scalar in model 2 to be the \(p\)-th power of that in
model 1. Two fixed scalar multiples of the same \(q\)-element
space cannot have this property: it would make Frobenius scalar
on all of \(\mathbf F_q\). This excludes \(r=2\).

Suppose \(r=3\), so \(A=pe\). Identify the residual space with
\(\mathbf F_q\) and decompose
\[
 \operatorname{End}_{\mathbf F_p}(\mathbf F_q)=Z\oplus U,\qquad
 Z=\{z\mapsto az\},\quad U=\{z\mapsto az^p\}.
\]
If the residual projective order is \(h>2\), tame conjugation
is trivial on \(Z\) and irreducible on \(U\), with multiplier
\(\lambda^{1-p}\) of order \(h\mid p+1\). Thus
\(W=Z_1\oplus U\) for a Cartan line \(Z_1\), and \(U\) has
one lower break. If \(h=2\), conjugation is \(+1\) on \(Z\)
and \(-1\) on \(U\), so either \(W=Z_1\oplus U\) or
\(W=Z\oplus U_1\). The projective order cannot be one,
since the residual representation is irreducible.

For every primitive class \(z\), use its fixed expansion
\[
 T_i(P)=d_i z\varpi^A+
       \sum_{A<j<qA}f_{i,j}(z)\varpi^j\pmod{\varpi^{qA}}.
 \tag{B-RD67.6}
\]
The coefficients below \(qA\) are independent of the representative,
because changing it adds a first-level point. No additivity
of the other coefficient functions is assumed. Every nonzero
matrix in \(Z\) or \(U\) is invertible, so its displacement
has valuation exactly \(qA\) for every \(z\ne0\).
For an element of break \(b<b_0=Ne\), the leading term first
changes in degree \(A+pb<qA\). Every other degree divisible
by \(p\) changes later. Cancellation forces the first degree
prime to \(p\) to be
\[
 j=A+(p-1)b,\qquad
 j f_{i,j}(z)=-\bar e\,d_i d_g^{p-1}z .
 \tag{B-RD67.7}
\]
In particular \(p\nmid j\). For \(b>b_0\), all degrees divisible
by \(p\) change after \(qA\), so that first prime-to-\(p\)
degree is \(j=qA-b\). Its coefficient is a nonzero multiple
of \(z\) for a Cartan element, and of \(z^p\) for an anti-Cartan
element. These statements refer to the same expansion (B-RD67.6).

First suppose \(U\subset W\) has one break \(b\). Its tame action
is nontrivial, so \(\eta^b\ne1\), whence \(b\ne b_0\).
If \(b<b_0\), fix \(z\ne0\) in (B-RD67.7). It requires
\(d_g^{p-1}\) to be constant on the \(q-1\) nonzero elements
of the leading coefficient space of \(U\), impossible.
If \(b>b_0\), the displacement equation at \(j=qA-b\) is
\[
 d_i^q\beta_i(g)z^p=j f_{i,j}(z)d_g,\qquad g\in U,
\]
where \(\beta_i:U\to\mathbf F_q\) is an additive isomorphism.
Fixing one \(g\ne0\), then varying \(g\), gives
\(\beta_i(g)=\kappa_i d_g\). The residual intertwiner requires
\(\beta_2(g)=a^{1-p}\beta_1(g)^p\), again making Frobenius
scalar on all of \(\mathbf F_q\). This excludes a single break.

If \(U\subset W\) has two breaks \(b_1<b_2\), necessarily \(h=2\);
both satisfy \(\eta^{b_j}=-1\), so neither equals \(b_0\).
If \(b_1<b_0\), (B-RD67.7) gives the first prime-to-\(p\)
term with coefficient proportional to \(z\). At \(b_2\) its
change is uniquely earlier than all terms divisible by \(p\).
It must therefore be the required first-level displacement,
but would be Cartan-linear in \(z\), contrary to its nonzero
anti-Cartan type. If \(b_1>b_0\), that first term and all
other terms change after \(qA\) at \(b_2\), also impossible.

It remains that \(h=2\) and \(W=Z\oplus U_1\).
Write \(b\) for the break of \(U_1\). Its tame weight is \(-1\),
while every Cartan break \(c\) has \(\eta^c=1\); hence
\(b\ne b_0\) and \(c\ne b\).
If \(b>b_0\), the first ordinary coefficient is proportional
to \(z^p\). A Cartan break \(c<b_0\) would make that same
coefficient proportional to \(z\). If \(b_0\le c<b\),
its change is uniquely before \(qA\); if \(c>b\), all changes
are after \(qA\). All possibilities are excluded, including
the boundary \(c=b_0\). Thus \(b<b_0\).
The first ordinary degree is now \(j=A+(p-1)b\), with coefficient
proportional to \(z\). A Cartan break below \(b\) would force
an earlier first ordinary term. Every Cartan break is therefore
greater than \(b\), and its unique first change forces
\[
 c=qA-j=p b_0-(p-1)b>b_0 .
 \tag{B-RD67.8}
\]
All nonzero Cartan elements have this same break. Their leading
displacement makes their scalar coefficients
\(\beta_i(g)=\kappa_i d_g\) on the full plane \(Z\).
Frobenius-semilinearity requires
\(\beta_2(g)=\beta_1(g)^p\), the same contradiction as before.
This exhausts every three-dimensional \(W\).
\end{proof}

\begin{lemma}[The final plane for the full wild image]
\label{lem:B-RD67-final-plane}
\status{Candidate; point fields and both ramification numberings explicit.}
If \(W=M_2(\mathbf F_p)\), then every primitive point field
\(L(P)/L\) is Galois of degree \(q\), with one lower and upper
break \(b_0=Ne\). The last nonzero upper group of \(F/L\) is
the full Cartan plane \(Z\). Its common lower break \(c\) satisfies
\(c>b_0\), and every anti-Cartan break is smaller than \(c\).
\end{lemma}
\begin{proof}
For a primitive residual vector \(v\ne0\), its stabilizer in
\(W\) is \(H_v=\{B:Bv=0\}\), a two-dimensional subspace.
Thus \(E=L(P)\) is Galois of degree \(q\). Its primitive
coordinate has valuation \(e\), prime to \(p\), and every
nonidentity automorphism displaces it by a first-level point
of valuation \(qe\). Fixed-coefficient uniformizer expansion
therefore gives the common lower break \(qe-e=b_0\).
A single lower break is also the single upper break.

Apply upper-numbering quotient compatibility
\cite[\S3.1, Proposition 3, p.~128]{Serre1961LocalClosed};
the full norm-polynomial and Herbrand proofs are also in
\cite[Proposition 3.7, Theorem 3.10 and Proposition 3.13,
pp.~14--16]{ColmezCorpsLocaux}.
The group \(S=W^{b_0}\) surjects under evaluation
\(B\mapsto Bv\) for every \(v\ne0\). There is no nonzero upper
group after \(b_0\), since \(\bigcap_v H_v=0\).

Over the good field \(L\), the determinant is cyclotomic by
the height-one determinant argument already proved in B18.
Its degree-\(p\) quotient is \(E_{\rm cyc}=L(\mu_{p^2})\).
This has degree \(p\), since \(\mu_p\subset L\) and \(p\nmid e_L\).
For \(x=\zeta_{p^2}-1\),
\[
 v_{E_{\rm cyc}}(x)=\frac{eN}{p-1},\qquad
 v_{E_{\rm cyc}}(gx-x)=\frac{peN}{p-1}\quad(g\ne1).
 \tag{B-RD67.9}
\]
The first integer is prime to \(p\), so this quotient also has
single break \(Ne=b_0\). Consequently trace is nonzero on \(S\).

At the last lower break \(c\), the leading ramification map
\(S\hookrightarrow k\) is additive and injective, with tame
weight \(\eta^c\). The nonzero invariant trace functional,
by averaging, supplies a nonzero tame-fixed vector of \(S\).
Hence \(\eta^c=1\), and injectivity makes all of \(S\) tame-fixed.
The centralizer of the irreducible residual tame element is
exactly \(Z\). Evaluation surjectivity forces dimension at least
two, so \(S=Z\). The Herbrand function has value \(b_0\) at
this last lower break. Since \(Z\) is a proper subgroup of \(W\),
its slope is strictly less than one on a nonempty interval before
\(c\), giving \(c>b_0\). This proves all assertions.
\end{proof}

\begin{lemma}[All two-dimensional anti-Cartan break patterns]
\label{lem:B-RD67-full-wild}
\status{Candidate; all three monomial slopes and coefficient functions retained.}
The full wild image \(W=M_2(\mathbf F_p)\) also cannot occur
in a different-special-character comparison.
\end{lemma}
\begin{proof}
Now \(A=qe\), with \(v_p(A)=2\), and the final Cartan plane
has the break \(c>b_0\) just proved. For a fixed primitive
coordinate \(z\ne0\), let \(J\) be the first nonzero degree
of exact \(p\)-valuation one, and \(K_0\) the first nonzero
degree prime to \(p\), in its expansion below \(qA\).
An absent degree is interpreted as infinity. All degrees
divisible by \(q\) change after \(qA\) at the break \(c\).
The \(q\)-element leading ramification space of \(Z\) cannot
be the root space of a nonzero polynomial of degree at most \(p\).
Applying this to a possible cancellation between the \(J\)
and \(K_0\) terms gives
\[
 J+pc\ge qA,\qquad K_0+c\ge qA,\qquad
 \min(J+pc,K_0+c)=qA .
 \tag{B-RD67.10}
\]
Every nonzero anti-Cartan matrix is invertible, so its displacement
also has exact valuation \(qA\).

First suppose the anti-Cartan plane has one break \(b<c\).
Its leading coefficient space \(D\) has dimension two over
\(\mathbf F_p\), but need not be a scalar multiple of \(\mathbf F_q\).
Its root polynomial is
\[
 P_D(X)=X^q+s_DX^p+t_DX,\qquad t_D\ne0.
 \tag{B-RD67.11}
\]
For example, with basis \(a,b\), put \(Y=X^p-a^{p-1}X\);
then \(P_D=Y^p-(b^p-a^{p-1}b)^{p-1}Y\).
The tame action on the anti-Cartan plane has no fixed vector,
so \(\eta^b\ne1\), and \(b\ne b_0\).
If \(b>b_0\), the leading degree \(A\) already changes after
\(qA\). An earlier cancellation between the \(J\) and \(K_0\)
terms would be a polynomial of degree at most \(p\) vanishing
on all of \(D\), impossible. Thus both of those changes are
at least \(qA\) at \(b\), contradicting (B-RD67.10) at \(c>b\).
Hence \(b<b_0\).

At this lower break, the degree-\(A\) change in degree \(A+qb\)
has nonzero coefficient \(\bar e\,d_i z d_g^q\).
No earlier change from the other two classes can cancel on
all of \(D\), by the same degree-\(p\) root bound.
Cancellation at \(A+qb\) must therefore be a multiple of
(B-RD67.11). Its nonzero linear coefficient forces
\[
 K_0=A+Nb,\qquad J\ge A+(q-p)b,
 \tag{B-RD67.12}
\]
and \(f_{i,K_0}(z)\) is one fixed nonzero multiple of \(z\).
The middle coefficient \(s_D\) is retained; equality in the
bound for \(J\) is allowed. At the later break \(c\),
\[
 J+pc-(K_0+c)\ge(p-1)(c-b)>0,\qquad
 A+qc-(K_0+c)=N(c-b)>0 .
 \tag{B-RD67.13}
\]
Thus the prime-to-\(p\) term is uniquely first and
\(K_0+c=qA\).

If the anti-Cartan plane has two breaks \(b_1<b_2<c\),
consider the three possible first-change degrees
\[
 A+qb,\qquad J+pb,\qquad K_0+b.
 \tag{B-RD67.14}
\]
At \(b_2\), (B-RD67.10) makes the minimum smaller than \(qA\);
it must be attained twice. It cannot involve the first term:
its larger slope would make it uniquely smallest at \(b_1\),
still below \(qA\). Hence
\[
 J+pb_2=K_0+b_2<A+qb_2.
 \tag{B-RD67.15}
\]
At \(b_1\), the middle term is strictly smaller than the last
and is below \(qA\). Cancellation forces equality with the first:
\[
 J=A+(q-p)b_1,\qquad
 K_0=J+(p-1)b_2,\qquad b_1<b_0.
 \tag{B-RD67.16}
\]
Choose fixed elements of the two breaks. The first cancellation
makes \(f_{i,J}(z)\) a fixed nonzero multiple of \(z\);
the second makes \(f_{i,K_0}(z)\) a fixed nonzero multiple
of that coefficient. These fixed elements work for every
\(z\ne0\); therefore the degrees and multiples are independent
of \(z\). At \(c\), the ordinary term is strictly first, since
\(J+pc-(K_0+c)=(p-1)(c-b_2)>0\); (B-RD67.10) again gives
\(K_0+c=qA\).

In either pattern, on the full final plane \(Z\) the displacement
equation is
\[
 d_i^q\beta_i(g)z=K_0 f_{i,K_0}(z)d_g .
 \tag{B-RD67.17}
\]
Thus \(\beta_i(g)=\kappa_i d_g\), and
\(\beta_i:Z\to\mathbf F_q\) is bijective.
The different-character intertwiner requires
\(\beta_2(g)=\beta_1(g)^p\), impossible for two fixed scalar
multiples of the same \(q\)-element space. This excludes \(r=4\).
\end{proof}

\begin{theorem}[Irreducible residual recovery for every prime-to-\(p\) index]
\label{thm:B-RD67}
\status{Candidate; full field, coordinate and ramification proof written.}
Let \(p\ge5\), let \(K\) be complete discretely valued of mixed
characteristic \((0,p)\) with perfect residue field, \(\mu_p\subset K\)
and \(p\nmid e_K\). Let \(M_1,M_2\) be height-two BT groups
which acquire one-dimensional supersingular good reduction.
If their \(p^2\)-torsion is isomorphic and the common residual
inertia is irreducible over \(\mathbf F_p\), then their finite
special inertia characters are equal.
\end{theorem}
\begin{proof}
Pass to the completed maximal unramified extension. This preserves
the absolute ramification and the given torsion comparison, while
the residue field becomes algebraically closed. The preceding
lemmas exclude each possible wild image size \(p^2,p^3,p^4\)
under the assumption that the special characters differ.
They exhaust the possible sizes by (B-RD67.1). All calculations
refer to the actual common division field and the given torsion
isomorphism, so they prove equality on inertia over the original
field as well.
\end{proof}

\begin{remark}[Scope and the role of the ramification index]
\label{rem:B-RD67-scope}
\status{Candidate scope; reducible residual comparisons are not settled here.}
The proof replaces the earlier minimal-index calculation by explicit
normalizations \(A=e p^{r-2}\), \(b_0=(p^2-1)e\), and
\(e_F=(p^2-1)p^2A\). The coefficient of the first changed leading
monomial is \(\bar e\), nonzero precisely because \(p\nmid e\).
No identification of lower and upper ramification numbers is made
before the point-field and cyclotomic quotient calculation.
The tame extension of degree \(p^2-1\) need not act faithfully on
the special groups. The theorem does not claim character recovery
for reducible residual inertia, ramification divisible by \(p\),
or general heights.
\end{remark}
